\documentclass[10pt,a4paper]{amsart}
\usepackage[margin=19mm]{geometry}
\usepackage{amsmath,amssymb,mathtools}
\usepackage[scr=rsfs,cal=euler]{mathalfa}
\usepackage{xfrac}
\usepackage[unicode,hidelinks,bookmarks=false]{hyperref}
\newcommand{\ie}{i.e.\ }
\newcommand{\cf}{cf.\ }
\newcommand{\resp}{respectively\ }
\newcommand{\Subsection}{\subsection}
\providecommand{\nobreakdash}{\nobreak\hbox{-}}
\setlength{\emergencystretch}{2em}
\multlinegap0pt
\usepackage{bm}
\usepackage{bbm}
\usepackage{dsfont}
\usepackage{xspace}
\usepackage{cancel}
\usepackage{mathtools}
\usepackage{cleveref}
\usepackage{amsthm}
\makeatletter
\@ifundefined{claim}{\newtheorem{claim}{Claim}}{}
\@ifundefined{theorem}{\newtheorem{theorem}{Theorem}}{}
\@ifundefined{lemma}{\newtheorem{lemma}[theorem]{Lemma}}{}
\makeatother
\makeatletter
\expandafter\ifx\csname proof\endcsname\relax
\renewenvironment{proof}[1][\proofname]{\par
	\pushQED{\qed}%
	\normalfont \topsep6\p@\@plus6\p@\relax
	\trivlist
	\item[\hskip\labelsep
	\bfseries
	#1\@addpunct{.}]\ignorespaces
}{%
	\popQED\endtrivlist\@endpefalse
}
\fi
\makeatother
\title[Brooks' theorem for signed graphs with $\Delta=3$]{Brooks' theorem for signed graphs with $\Delta=3$}
\author{Reza Naserasr, Huan Zhou}
\date{Source: arXiv:2510.16490v1 (CC BY 4.0)}
\begin{document}
\maketitle

\noindent\emph{Source and licence.} Brooks' theorem for signed graphs with $\Delta=3$. Version arXiv:2510.16490v1 (CC BY 4.0), distributed under the Creative Commons Attribution 4.0 International licence, as recorded in the arXiv licence metadata for this exact version and shown on the abstract page https://arxiv.org/abs/2510.16490v1. Full rights notice: licenses/arxiv-2510.16490-CC-BY-4.0.txt.\par\smallskip

\noindent\textit{Editorial scope.} Counted unit: the Lemma whose source label is \texttt{lem:key} in the article (source0.tex lines 854--862, source label \texttt{lem:key}), delivered together with the article's own complete original proof of it (source0.tex lines 864--934, one \texttt{proof} environment of 71 lines). No mathematics is written, repaired or re-derived: statement and proof are the article's own text line for line. Required context retained uncounted: lem:T (lines 312--314); lem:3-colorable (lines 389--391); thm:main (lines 844--846). Every definition, display and label of those lines is retained unchanged. Omitted from this package and not redistributed: 1--311, 315--388, 392--843, 847--853, 863--863, 935--1626. Nothing inside a retained excerpt is omitted or altered: every retained body is delivered exactly as the article's lines are written, so no omission receipt of any kind accompanies this package. Citations: the retained excerpts carry no citation command at all, so no bibliography entry of the article is transcribed and no reference is redistributed. Reference adaptation: none was needed and none was made; every reference inside the delivered unit resolves inside it, so no unresolved reference or citation is produced. Printed slips kept literal: no printed word, symbol, hypothesis or number of the counted statement or of its proof was altered. Layout and class: the article's own class is \texttt{article} with packages including mathrsfs, amsmath, amsthm, amssymb, bm, mathtools, thm-restate, authblk, graphicx, xcolor, tikz, caption, subcaption, array; the wrapper is the lane's pinned \texttt{amsart} editorial wrapper at 10pt on A4, which restates the theorem-like environments the retained text needs and adds the article's own macro definitions that the retained text needs (none), with extra wrapper packages (bm, bbm, dsfont, xspace, cancel, mathtools, cleveref). The delivered numbering is the unit's own. Assets: no retained span contains a figure environment, a graphics-inclusion command, a PDF-page inclusion, a table or a TikZ picture; no image file is copied into this package, the package declares no asset dependency and no image file is redistributed with it. Licence: version arXiv:2510.16490v1 of this article is distributed under the Creative Commons Attribution 4.0 International licence, as recorded in the arXiv licence metadata for this exact version and shown on the abstract page https://arxiv.org/abs/2510.16490v1. A full rights notice is delivered at licenses/arxiv-2510.16490-CC-BY-4.0.txt. Characters: every retained excerpt is pure ASCII, so the delivered engine, XeLaTeX with its default Latin Modern text font, reports no missing character.
	\begin{lemma}\label{lem:T}
		Every homomorphism of $\widehat{T}$ to $\widehat{K}^s_{10;3}$ is surjective.  
	\end{lemma}

	\begin{lemma}\label{lem:3-colorable}
		Assume $\widehat{G}$ is a signed simple graph such that $|V(\widehat{G})|\leq 5$,  $|E(\widehat{G})|\leq 7$, $\widehat{G}\ncong\widehat{T}$, and $\widehat{G}$ contains no subgraph isomorphic to $(K_4,-)$. Then $\widehat{G}$ is circular 3-colorable. 
	\end{lemma}

\begin{theorem}\label{thm:main}
	If $\widehat{G}$ is a 10/3-critical signed graph which is isomorphic neither to $(K_4,-)$ nor to $\widetilde{K}_2$, then $$p(\widehat{G})=3|V(\widehat{G})|-2|E(\widehat{G})|\leq -1.$$
\end{theorem}

\begin{lemma}\label{lem:key}
	Assume $\widehat{H}\subseteq \widehat{G}$. Then 
	\begin{itemize}
		\item[(I)] $p(\widehat{H})\ge 0$, if $\widehat{H} \cong \widehat{G}$;
		\item[(II)] $p(\widehat{H})\ge 1$, if $\widehat{G}\in P_2(\widehat{H})$ or 
		$\widehat{H}\cong \widehat{T}$;
		\item[(III)] $p(\widehat{H})\ge 2$, otherwise.
	\end{itemize}
\end{lemma}

\begin{proof}
	The statement (I) is the assumption on $\widehat{G}$.
	If $\widehat{G}\in P_2(\widehat{H})$, then $p(\widehat{H})=p(\widehat{G})-(3\times 1-2\times 2)=p(\widehat{G})+1\ge 1$. The potential value of $\widehat{T}$ is $3\times 5-2\times 7=1$.
	
	
	Suppose (III) is not true. Let $\widehat{H}$ be a maximal proper subgraph of $\widehat{G}$ satisfying $\widehat{G} \not\in P_2(\widehat{H})$, $\widehat{H}\ncong \widehat{T}$, and $p(\widehat{H})\leq 1$. Since adding an edge decreases the value of the potential, we may assume $\widehat{H}$ is an induced subgraph of $\widehat{G}$. 
	
	
	Observe that $p(\widehat{K}_1)=3$, $p(\widehat{K}_2)=4$, $p(\widehat{K}_3)=3$. Assume $\widehat{H}$ has $4$ or $5$ vertices. Since $p(\widehat{H})\leq 1$, we have $E(\widehat{H})\leq 7$ and hence  by \Cref{lem:3-colorable}, $\widehat{H}$ maps to a negative triangle. If it has more than $6$ vertices, then since it is a proper subgraph, it maps to $K_{10;3}^{s}$. In both cases, we conclude that there exists a homomorphism $\phi$ of $\widehat{H}$ to $\widehat{K}^s_{10;3}$ which identifies at least two vertices. In other words, $|\phi(\widehat{H})|<|V(\widehat{H})|$.
	We extend $\phi$ to $\widehat{G}$ by taking it as an identity mapping on the rest of $\widehat{G}$.
	
	
	
	Define $\widehat{G}^*$ to be the signed graph obtained from $\phi(\widehat{G})$ as follows: 
	\begin{itemize}
		\item if $\phi(\widehat{H})$ contains no subgraph isomorphic to $\widehat{T}$, then $\widehat{G}^*=\phi(\widehat{G})$;
		\item if $\widehat{T}_{1}\subseteq \phi(\widehat{H})$ and $\widehat{T}_1\cong \widehat{T}$, then $\widehat{G}^*=\phi(\widehat{G})-E_d$, where $E_d=E(\phi(\widehat{H}))-E(\widehat{T}_{1})$.
	\end{itemize}
	
	In other words, after identifying,  the $\widehat{H}$-part of $\widehat{G}$ with a subgraph of  $\widehat{K}^{s}_{10;3}$, if the image of $\widehat{H}$ contains more edges than $\widehat{T}$, then we delete those extra edges.
	
	
	
	By \Cref{lem:T}, if $\widehat{G}^*\to \widehat{K}^s_{10;3}$, then $\phi(\widehat{G})\to \widehat{K}^s_{10;3}$, which would imply that $\widehat{G}$ is $10/3$-colorable. Hence $\widehat{G}^*\not\to \widehat{K}^s_{10;3}$. 
	
	
	Since  $\widehat{G}^*\not\to \widehat{K}^s_{10;3}$, $\widehat{G}^*$ contains a $10/3$-critical graph $\widehat{G}'$. As $|V(\widehat{G}^*)|<|V(\widehat{G})|$, we have $|V(\widehat{G}')|<|V(\widehat{G})|$. 
	
	
	
	\begin{claim}\label{clm:digon}
		$\widehat{G}'\ncong \widetilde{K}_2$.
	\end{claim}
	
	\noindent\emph{Proof of the claim.} Assume to the contrary, and suppose vertices $u, v$ induce a digon in $\widehat{G}^*$. Since there is no digon in $\widehat{G}$ and $\widehat{K}^s_{10;3}$, one of $u$ and $v$, say $u$, is in $\phi(\widehat{H})$, the other one is in $V(\widehat{G})-V(\widehat{H})$. Moreover, $v$ is adjacent to at least two vertices of $\widehat{H}$. So $$p(\widehat{G}[V(\widehat{H})+v])\leq p(\widehat{H})+3\times 1-2 \times 2=p(\widehat{H})-1<p(\widehat{H}) .$$
	
	By the choice of $\widehat{H}$,  either $\widehat{G}[V(\widehat{H})+v]$ is isomorphic to $\widehat{G}$ or $\widehat{G}\in P_2(\widehat{G}[V(\widehat{H})+v])$. The latter is impossible because $p(\widehat{G}[V(\widehat{H})+v])\leq p(\widehat{H})-1\leq 0$. Hence $\widehat{G}[V(\widehat{H})+v]=\widehat{G}$. Moreover, $d(v)=2$, as otherwise $p(\widehat{G}[V(\widehat{H})+v])\leq p(\widehat{H})-3\leq -2$. That means $\widehat{G}\in P_2(\widehat{H})$, which contradicts the assumption.  
	
	
	\begin{claim}\label{clm:K_4}
		$\widehat{G}'\ncong (K_4,-)$.
	\end{claim}
	
	\noindent\emph{Proof of the claim.} Assume to the contrary that $\widehat{K}$ is a subgraph of  $\widehat{G}^*$ isomorphic to $(K_4,-)$. Since neither $\widehat{G}$ nor $\widehat{K}^s_{10;3}$ contains a subgraph isomorphic to $(K_4,-)$, we have $1\leq |V(\widehat{H}) \cap V(\widehat{K})|\leq 3$. Based on the size of $V(\widehat{H}) \cap V(\widehat{K})$, we consider three cases.
	\begin{itemize}
		\item[(1)] $|V(\widehat{H}) \cap V(\widehat{K})|=1$.
		
		In this case, $\widehat{K}\cap (\widehat{G}-\widehat{H})$  is a (negative) triangle, and hence there are at least three edges connecting $\widehat{H}$ and the triangle. Thus 
		$$p(\widehat{G}[V(H)\cup V(K_4)])\leq p(\widehat{H})+3\times 3-2\times 6=p(\widehat{H})-3\leq -2.$$ 
		
		
		\item[(2)] $|V(\widehat{H}) \cap V(\widehat{K})|=2$.
		
		The intersection $\widehat{K}\cap (\widehat{G}-\widehat{H})$ induces a $K_{2}$ which is connected to $\widehat{H}$ by at least four edges. Hence $$p(\widehat{G}[V(H)\cup V(K_4)])\leq p(\widehat{H})+3\times 2-2\times 5=p(\widehat{H})-4\leq -3.$$ 
		
		
		\item[(3)] $|V(\widehat{H}) \cap V(\widehat{K})|=3$.
		
		Similarly, noting that the  intersection $\widehat{K}\cap (\widehat{G}-\widehat{H})$ is a vertex connected to $\widehat{H}$ with at least three edges, we have
		$$p(\widehat{G}[V(H)\cup V(K_4)])\leq p(\widehat{H})+3\times 1-2\times 3=p(\widehat{H})-3\leq -2.$$ 
		
	\end{itemize}
	
	In all cases, we find a subgraph including $\widehat{H}$ with potential value at most $-2$. This either contradicts \Cref{lem:key}(I), or  \Cref{lem:key}(II), or the choice of $\widehat{H}$, hence completing the proof of Claim~\ref{clm:K_4}.
	\bigskip
	
	It follows from \Cref{clm:digon} and \Cref{clm:K_4} that the $10/3$-critical graph $\widehat{G}'$ satisfies all the conditions of Theorem~\ref{thm:main}. Since $|V(\widehat{G}')|<|V(\widehat{G})|$, we have $p(\widehat{G}')\leq -1$.
	
Since potential is a linear function, we have 	$$p(\widehat{G}[V(\widehat{G}'-\phi(\widehat{H}))+V(\widehat{H})])\leq p(\widehat{G}')-p(\widehat{G}'\cap \phi(\widehat{H}))+p(\widehat{H})\leq p(\widehat{G}')-1+p(\widehat{H})\leq p(\widehat{H})-2\leq -1, $$
	which again either contradicts  \Cref{lem:key}(I), or  \Cref{lem:key}(II), or the choice of $\widehat{H}$. This completes the proof.
\end{proof}

\end{document}
